\section*{ABSTRACT}
\phantomsection\addcontentsline{toc}{section}{\numberline{}ABSTRACT}

This report presents a comparative study of Orthogonal Frequency Division
Multiplexing (OFDM) and Single-Carrier Frequency Division Multiple Access
(SC-FDMA), two important waveform technologies used in modern wireless
communication systems. OFDM is widely recognized for its high spectral
efficiency, robustness against multipath fading, and simple frequency-domain
equalization. However, one major drawback of OFDM is its high
Peak-to-Average Power Ratio (PAPR), which reduces power amplifier
efficiency and may introduce nonlinear distortion.

SC-FDMA was developed to overcome this limitation by applying a Discrete
Fourier Transform (DFT) precoding stage before OFDM modulation. As a
result, SC-FDMA preserves many advantages of OFDM while achieving lower
PAPR and improved transmitter power efficiency. This characteristic makes
SC-FDMA especially suitable for uplink transmission in battery-powered
mobile devices.

The report reviews the operating principles of both systems and compares
their performance using two common evaluation metrics: Bit Error Rate
(BER) and PAPR. A MATLAB-based simulation framework is proposed using
16-QAM modulation under additive white Gaussian noise conditions. Expected
results indicate that OFDM and SC-FDMA provide similar BER performance
under identical channel assumptions, while SC-FDMA demonstrates clear
advantages in PAPR reduction.

The study highlights the engineering tradeoff between transmission
efficiency and implementation structure, and explains why OFDM is commonly
used in downlink systems whereas SC-FDMA was adopted for LTE uplink.

\newpage
\clearpage